\scanpage{18}
\diagram{\begin{tikzcd}[ampersand replacement=\&,column sep=small,row sep=small]
1 \arrow[r] \& N \arrow[r] \& \overset{\deleted{\unsure{\sim}}}{G}_0 \arrow[r] \& \unsure{\mathfrak S|\mathfrak G} \arrow[r] \& 1\\
\& \& \& \pi \arrow[u] \arrow[ul,dashed,"\varphi"'] \&
\end{tikzcd}}
invariant par $\widetilde G$ [opérant sur $\pi$ (via $\widetilde H$ et $\widetilde G\to\widetilde H$) et \deleted{\unsure{$N$}} \inserted{$G_0$} (via $G$ \inserted{et $\widetilde G\to G_0$})]
% slanted
\margin{si $\pi$ \unsure{comm.|connexe} et discret, l'opération \inserted{de $\unsure{G|\widetilde G}$} sur $\pi$ se fait via $\unsure{H|\widetilde H}$, donc via $\pi_0(H)$ (ou $\pi_0(\widetilde H)$) et celle de $\unsure{G|\widetilde G|\mathfrak G}$ sur $\mathfrak z(N)$ se fait via $\mathfrak S$ (\unsure{car} $N$ \deleted{\unsure{agit}} \unsure{opère} triv$^{\rm t}$) et si $\mathfrak S$ discret et $\widetilde H_0$ connexe, on a $\pi_0(H)=\pi_0(\widetilde H)=\mathfrak S$) donc la condition est que $\varphi:\pi\to\mathfrak z(N)$ soit $\mathfrak S$-hom.}

Cas particulier : $\pi\subset\widetilde H_0$, i.e. $\pi\to\mathfrak S$ est trivial, alors la donnée de $\varphi$ équivaut à la donnée d'un hom
\[
\varphi:\pi\to\mathfrak z(N),
\]
\noindent\rule{\linewidth}{.4pt}

\diagram{\scriptsize\begin{tikzcd}[ampersand replacement=\&,column sep=small,row sep=small,/tikz/execute at end picture={\draw[thin] (current bounding box.north west) -- (current bounding box.south east);}]
\& \& \& \deleted{\widetilde H_0} \arrow[dl,hook] \arrow[dr] \& \&\\
1 \arrow[r] \& \pi \arrow[r] \& \widetilde H \arrow[rr] \& \& H \arrow[r] \& 1\\
1 \arrow[r] \& \pi \arrow[u,equal] \arrow[r] \& \widetilde H' \arrow[u,"\widetilde f"] \arrow[rr] \& \& H' \arrow[u,"\deleted{\unsure{f}}", "f"'] \arrow[r] \& 1\\
\& \& \& \widetilde H'_0=f^{-1}(\widetilde H_0) \arrow[uuu] \arrow[ul,hook] \arrow[ur] \& \&
\end{tikzcd}}
\[
\deleted{\unsure{\mathrm{int}(\widetilde g)}\,\varphi(\pi)=\unsure{-}}
\]
\diagram{\scriptsize\begin{tikzcd}[ampersand replacement=\&,row sep=small,column sep=small,/tikz/execute at end picture={\draw[thin] (current bounding box.north west) -- (current bounding box.south east);}]
\& \& 1 \& \&\\
1 \arrow[r] \& \pi \arrow[r] \arrow[dr,"\varphi"'] \& \deleted{\widetilde H} \arrow[u] \arrow[r] \& H \arrow[r] \& 1\\
1 \arrow[r] \& \deleted{\unsure{\pi}} \arrow[u] \arrow[r] \arrow[d] \& \widetilde G \arrow[u] \arrow[r] \& \unsure{G?|G} \arrow[u] \arrow[r] \& 1\\
\& \deleted{\unsure{N=}} \& N \arrow[u] \arrow[r] \& N \arrow[u] \&\\
\& \& 1 \arrow[u] \& \&
\end{tikzcd}}

\diagram{\scriptsize\begin{tikzcd}[ampersand replacement=\&,row sep=small,column sep=small,/tikz/execute at end picture={\draw[thin] (current bounding box.north west) -- (current bounding box.south east);}]
1 \arrow[r] \& \mathscr H_1/\mathscr H_1\cap\mathscr H_2 \arrow[r] \& \overset{\unsure{p_2|j_2}}{\mathscr H/\mathscr H_2} \arrow[r] \& \overset{\unsure{p_1|j_1}}{\mathscr H/\mathscr H_1\mathscr H_2} \arrow[r] \& 1\\
1 \arrow[r] \& \mathscr H_1 \arrow[r,hook] \& \deleted{\mathscr H} \arrow[u] \arrow[r] \& \mathscr H/\mathscr H_1 \arrow[u] \arrow[r] \& 1\\
1 \arrow[r] \& \mathscr H_1\cap\mathscr H_2 \arrow[u] \arrow[r] \& \mathscr H_2 \arrow[u,hook] \arrow[r] \& \mathscr H_2/\mathscr H_1\cap\mathscr H_2 \arrow[u] \arrow[r] \& 1\\
\& 1 \arrow[u] \& 1 \arrow[u] \& 1 \arrow[u] \&
\end{tikzcd}}

\diagram{\scriptsize\begin{tikzcd}[ampersand replacement=\&,row sep=tiny,column sep=small,/tikz/execute at end picture={\draw[thin] (current bounding box.north west) -- (current bounding box.south east);}]
\& \& 1 \& 1 \&\\
\& \& \mathfrak S \arrow[u] \arrow[r,equal] \& \mathfrak S \arrow[u] \&\\
1 \arrow[r] \& \pi \arrow[r] \& \deleted{\widetilde H} \arrow[u] \arrow[r] \& H \arrow[u] \arrow[r] \& 1\\
1 \arrow[r] \& \pi \arrow[u,equal] \arrow[r] \& \widetilde H_0 \arrow[u] \arrow[r] \& H_0 \arrow[u] \arrow[r] \& 1\\
\& \& 1 \arrow[u] \& 1 \arrow[u] \&
\end{tikzcd}}
\diagram{\scriptsize\begin{tikzcd}[ampersand replacement=\&,column sep=1em]
\mathrm{Ext}(\mathfrak S,N) \arrow[d] \&[-2em] {}\times\mathrm{Hom}_{\mathfrak S}(\pi,\mathfrak z(N)) \arrow[r,"\sim"] \& \deleted{\mathrm{Ext}(H,\mathrm{mod}\,\widetilde H_0,N)} \arrow[d]\\
\substack{\text{Esp. homogène}\\\text{sur }H^2(\mathfrak S,N)} \& \& \mathrm{Ext}(H,N)\\
\& \& \substack{\text{espace homogène}\\\text{sur }H^2(H,\mathfrak z(N))}
\end{tikzcd}}
\diagram{\scriptsize\begin{tikzcd}[ampersand replacement=\&,row sep=tiny,column sep=tiny]
\& \& \deleted{\unsure{\mathscr H}} \& \cdot \&\\
1 \arrow[r] \& N \arrow[r] \& G \arrow[r] \arrow[d] \& \widetilde H \arrow[u] \arrow[r] \& 1\\
\& \& H \& \overset{\deleted{\unsure{\sim}}}{H} \&
\end{tikzcd}}
\[
\deleted{H^2(\mathfrak S,\mathfrak z N)\to H^{2}(H,\mathfrak z N)}
\]
car $H\to\mathfrak S$
\[
\deleted{\mathrm{Ext}(\mathfrak S,N)\times\mathrm{Hom}(\pi,N)\to H^2(H,\mathfrak z(N))}.
\]
% archive sheet number: 17
